% !TEX encoding = UTF-8 Unicode

%%%%%%%%%%%%%%%%%%%%%%%%%%%%%%%%%%%%%%%%%%%%%%%%%%%%%%%%%%%%%%%
%                                                             %
% MINIATURAS MATEM\'ATICAS DE LA GACETA DE LA RSME            %
%                                                             %
%%%%%%%%%%%%%%%%%%%%%%%%%%%%%%%%%%%%%%%%%%%%%%%%%%%%%%%%%%%%%%%

\documentclass[autocontact]{gaceta}

%---------------------
% La Gaceta se har\'a con esto activado, 
% pero no hace falta ahora si no quieres:
\usepackage[T1]{fontenc}
\usepackage{lmodern}

% Puedes usar esto (u otra codificaci\'on) si quieres:
\usepackage[utf8]{inputenc}
\usepackage[spanish]{babel}
%\decimalpoint
%---------------------

%---------------------
% Esto ya lo carga el estilo de La Gaceta:
%\usepackage{amsmath, amsthm, amssymb}
%\usepackage{url} % <-- Para p\'aginas web o similar: \url{...}
%\usepackage{graphicx}
%---------------------

%---------------------
% Carga esto para gr\'aficos si lo necesitas:
\usepackage{wrapfig}
%---------------------

%---------------------
% Carga esto para gr\'aficos si lo necesitas:
\usepackage{tikz}
%\usepackage{pgfplots}
%---------------------

%---------------------
% <<< ESTO SE AJUSTAR\'A AL EDITAR CADA N\'UMERO DE LA GACETA >>>
%\provisional
\newcounter{pageforminiatura}
\setcounter{pageforminiatura}{326} % Siempre ser\'a p\'agina par
\setcounter{page}{1} % No se usa (dejarlo como 1)
\journame{La Gaceta de la RSME}
\yearofpublication{2022}
\volume{25}
\issuenumber{2}
%---------------------

%---------------------
% Deja esto as\'{\i}:
%\belongstopart{Art\'{\i}culos} 
%\belongstopart{Secciones} 
%\belongstopart{Noticias de la Sociedad} 
%\belongstopart{Actualidad} 
\nopart
%---------------------

%---------------------
% Deja esto as\'{\i}: 
\supertitle{\large Miniaturas matem\'aticas de \thejourname\\[1mm]
\hfill\footnotesize
Vol.~\thevolume\ (\theyearofpublication), 
N\'um.~\theissuenumber, P\'ag.~\thepageforminiatura\theprovisional}
%
\noeditor
%---------------------

%---------------------
\title{La razón áurea y las medias metálicas\\en el salino de Arquímedes}
%\author{} % Poner \author{} o no poner nada, da igual.
% Para ahorrar espacio, el autor solo se pone en la firma:
\contact{Juan Luis Varona, Dpto. de Matemáticas y Computación, Universidad de La Rioja}
{jvarona@unirioja.es}
{https://www.unirioja.es/cu/jvarona/}
%---------------------

%%%%%%%%%%%%%%%%%%%%%%%%%%%%%%%%%%%%%%%%%%%%%%%%%%%%%%%%%%%%%%%%%%%%%%%
\begin{document}
%%%%%%%%%%%%%%%%%%%%%%%%%%%%%%%%%%%%%%%%%%%%%%%%%%%%%%%%%%%%%%%%%%%%%%%

%%%%%%%%%%%%%%%%%%%%%%%%%%%%%%%%%%%%%%%%%%%%%%%%%%%%%%%%%%%%%%%%%%%%%%%
% Deja esto as\'{\i}: 
\previousvspace{-1.8cm}
\makesupertitle
\previousvspace{-0.25cm}
\maketitle
\vspace*{-0.4cm}
\thispagestyle{empty} % <<<<< Importante
%%%%%%%%%%%%%%%%%%%%%%%%%%%%%%%%%%%%%%%%%%%%%%%%%%%%%%%%%%%%%%%%%%%%%%% 

%% NOTA: realmente, que $k$ sea entero no tiene importancia
%Para $k$ entero positivo, sean $\varphi_k$ y $\phi_k$
%las dos raíces $({k\pm \sqrt{k^2+4}})/{2}$ de la
%ecuación $x^2-kx-1 = 0$ (con $\varphi_k$ la positiva).
%Con esta notación, $\varphi_k$ son las denominadas medias metálicas; % (véase~\cite{P})
%cuando $k=1$, $\varphi_1$ es la razón áurea.
%%% ALTERNATIVA:
%% cuando $k=1$, $\varphi_1 = (1+\sqrt{5})/2$ es la razón áurea.

% NOTA: realmente, que $k$ sea entero no tiene importancia
Para $k$ entero positivo, sea $\varphi_k = ({k+\sqrt{k^2+4}})/{2}$ 
la raíz positiva de la ecuación $x^2-kx-1 = 0$.
Estas $\varphi_k$ satisfacen 
% $\varphi_k/1 = 1/(\varphi_k-k)$ 
% $(\varphi_k-k)\varphi_k = 1$ 
$\varphi_k = k+1/\varphi_k$ % <--- De aquí sale la fracción continua (k;k,k...)
y son las denominadas medias metálicas; 
% (véase~\cite{P})
cuando $k=1$, $\varphi_1$ es la razón áurea.

%% ASÍ ES COMO LO PUSE EN EL DE PLAZA:
%Sean $\varphi_k$ y $\phi_k$, respectivamente, las raíces positiva y negativa 
%de la ecuación $x^2=kx+1$. Con esta notación,
%y si $k$ es un entero positivo, $\varphi_k$ son las denominadas medias metálicas;
%% NOTA: realmente, que $k$ sea entero no tiene importancia
%cuando $k=1$, $\varphi_1$ es la razón áurea.

\begin{wrapfigure}{r}{0.20\textwidth}%
\vspace*{-1mm}%
\hspace*{-2mm}%
\newcommand{\kmedios}{1.5}% % k = 3
\newcommand{\mediaMetalica}{3.3027756}% % (k+sqrt(4+k^2))/2
\newcommand{\diametroEntrante}{0.3027756}% % (-k+sqrt(4+k^2))/2 = \mediaMetalica - k = 1/\mediaMetalica
\newcommand{\radio}{1.80278}% \kmedios + \diametroEntrante = sqrt(4+k^2)/2 
\begin{tikzpicture}[scale=0.66]
\footnotesize % <--- Tiene influencia sobre el texto de los "node"
\draw[thick] (0,0) -- (\kmedios,0) node[below]{$O$} -- (0,1) node[above left=-2]{$P$} -- cycle;
% \draw[green] (\kmedios,0) circle(\radio); % <- La circunferencia entera
\draw (\radio+\kmedios,0) arc (0:180:\radio);
\draw (-\radio+\kmedios,0) -- (\radio+\kmedios,0);
\draw (0,0.4) node[right]{$1$};
\draw (0.5*\kmedios,0) node[below]{$k/2$\,};
\end{tikzpicture}
\end{wrapfigure}
Sobre un triángulo rectángulo de catetos $1$ y $k/2$ se construye una semicircunferencia como muestra el gráfico adjunto.
%
%
A partir de esa semicircunferencia construimos un salino de Arquímedes 
~---\'Oscar Ciaurri nos lo presenta en su artículo «El ``Libro de los lemas'': un ejercicio de visualización»,
\textit{La Gaceta de la RSME} \textbf{17} (2014), 221--245---
con lóbulo central de diámetro~$k$,
% (véase~\cite{C}), 
y prolongamos el radio $PO$ (hipotenusa del triángulo inicial) hasta que corta al borde del lóbulo central en el punto~$Q$. 
En esas circunstancias, la longitud de $PQ$ es~$\varphi_k$.
Aquí mostramos los salinos correspondientes a $k=1$, $2$ y~$3$:
\begin{center}
\smallskip
\newcommand{\arito}{1.2pt}
\begin{tikzpicture}[scale=1.27]
\footnotesize % <--- Tiene influencia sobre el texto de los "node"
\begin{scope}[xshift=-2.5cm]
\newcommand{\kmedios}{0.5} % k = 1
\newcommand{\mediaMetalica}{1.61803} % (k+sqrt(4+k^2))/2
\newcommand{\diametroEntrante}{0.61803} % % (-k+sqrt(4+k^2))/2 = \mediaMetalica - k = 1/\mediaMetalica
\newcommand{\radio}{1.11803} % \kmedios + \diametroEntrante = sqrt(4+k^2)/2 
\newcommand{\angulo}{63.4349} % atan(1/(k/2)) = atan(2/k) = 1.10715 = 63.4349 grados
\newcommand{\qx}{0.447214} % cos(\angulo)
\newcommand{\qy}{0.894427} % sin(\angulo)
\coordinate (P) at (0,1);
\coordinate (Q) at (\kmedios+\kmedios*\qx,-\kmedios*\qy);
\draw[dashed] (0,0) -- (P) node[above left=-2]{$P$}; 
\draw (0,1) -- (\kmedios,0) node[below]{$O$\,\,\,\,} -- (Q) node[below right=-1]{$Q$};
\draw[dashed] (-\radio+\kmedios,0) -- (\radio+\kmedios,0);
\draw[thick] (\radio+\kmedios,0) arc (0:180:\radio);
\draw[thick] (\kmedios+\kmedios,0) arc (0:-180:\kmedios); % lóbulo central
\draw[thick] (0,0) arc (0:180:0.5*\diametroEntrante); % lóbulo izquierdo
\draw[thick] (\kmedios+\kmedios+\diametroEntrante,0) arc (0:180:0.5*\diametroEntrante); % lóbulo derecho
%\draw (0,0.4) node[right]{$1$};
%\draw (0.5*\kmedios,0) node[below]{$k/2$};
%\draw (\kmedios,0) -- (Q) node[below right=-1]{$Q$};
\draw[fill=white] (P) circle (\arito);
\draw[fill=white] (Q) circle (\arito);
\draw (0.55,0.4) node{$\varphi_k$};
\end{scope}
\begin{scope}[xshift=0cm]
\newcommand{\kmedios}{1.0} % k = 2
\newcommand{\mediaMetalica}{2.41421356} % (k+sqrt(4+k^2))/2
\newcommand{\diametroEntrante}{0.41421356} % % (-k+sqrt(4+k^2))/2 = \mediaMetalica - k = 1/\mediaMetalica
\newcommand{\radio}{1.41421356} % \kmedios + \diametroEntrante = sqrt(4+k^2)/2 
\newcommand{\angulo}{45.0} % atan(1/(k/2)) = atan(2/k) = pi/4 = 45.0 grados
\newcommand{\qx}{0.707107} % cos(\angulo)
\newcommand{\qy}{0.707107} % sin(\angulo)
\coordinate (P) at (0,1);
\coordinate (Q) at (\kmedios+\kmedios*\qx,-\kmedios*\qy);
\draw[dashed] (0,0) -- (P) node[above left=-2]{$P$}; 
\draw (0,1) -- (\kmedios,0) node[below]{$O$\,\,} -- (Q) node[below right=-1]{$Q$};
\draw[dashed] (-\radio+\kmedios,0) -- (\radio+\kmedios,0);
\draw[thick] (\radio+\kmedios,0) arc (0:180:\radio);
\draw[thick] (\kmedios+\kmedios,0) arc (0:-180:\kmedios); % lóbulo central
\draw[thick] (0,0) arc (0:180:0.5*\diametroEntrante); % lóbulo izquierdo
\draw[thick] (\kmedios+\kmedios+\diametroEntrante,0) arc (0:180:0.5*\diametroEntrante); % lóbulo derecho
%\draw (0,0.4) node[right]{$1$};
%\draw (0.5*\kmedios,0) node[below]{$k/2$};
%\draw (\kmedios,0) -- (Q) node[below right=-1]{$Q$};
\draw[fill=white] (P) circle (\arito);
\draw[fill=white] (Q) circle (\arito);
\draw (0.95,0.35) node{$\varphi_k$};
\end{scope}
\begin{scope}[xshift=3.2cm]
\newcommand{\kmedios}{1.5} % k = 3
\newcommand{\mediaMetalica}{3.3027756} % (k+sqrt(4+k^2))/2
\newcommand{\diametroEntrante}{0.3027756} % % (-k+sqrt(4+k^2))/2 = \mediaMetalica - k = 1/\mediaMetalica
\newcommand{\radio}{1.80278} % \kmedios + \diametroEntrante = sqrt(4+k^2)/2 
\newcommand{\angulo}{33.6901} % atan(1/(k/2)) = atan(2/k) = 0.588003 = 33.6901 grados
\newcommand{\qx}{0.83205} % cos(\angulo)
\newcommand{\qy}{0.55470} % sin(\angulo)
\coordinate (P) at (0,1);
\coordinate (Q) at (\kmedios+\kmedios*\qx,-\kmedios*\qy);
\draw[dashed] (0,0) -- (0,1) node[above left=-2]{$P$}; 
\draw (0,1) -- (\kmedios,0) node[below]{$O$\,} -- (Q) node[below right=-1]{$Q$};
\draw[dashed] (-\radio+\kmedios,0) -- (\radio+\kmedios,0);
\draw[thick] (\radio+\kmedios,0) arc (0:180:\radio);
\draw[thick] (\kmedios+\kmedios,0) arc (0:-180:\kmedios); % lóbulo central
\draw[thick] (0,0) arc (0:180:0.5*\diametroEntrante); % lóbulo izquierdo
\draw[thick] (\kmedios+\kmedios+\diametroEntrante,0) arc (0:180:0.5*\diametroEntrante); % lóbulo derecho
\draw (0,0.4) node[right]{$1$};
\draw (0.5*\kmedios,0) node[below]{$k/2$\,};
\draw (1.5*\kmedios,0) node[below]{\,$k/2$};
\draw (-0.5*\diametroEntrante,0) node[below]{$a$};
\draw (2.0*\kmedios+0.5*\diametroEntrante,0) node[below]{$a$};
%\draw (\kmedios,0) -- (Q) node[below right=-1]{$Q$};
\draw (-\diametroEntrante,0) node[below left=-1pt]{$A$\!};
\draw (0,0) node[below right=-0.5pt]{$B$};
\draw (2.0*\kmedios,0) node[below left=-0.5pt]{$C$};
\draw (2.0*\kmedios+\diametroEntrante,0) node[below right=-1pt]{\!$D$};
\draw[dotted] (-\diametroEntrante,0) -- (0,1) -- (2.0*\kmedios+\diametroEntrante,0);
\draw[fill=white] (P) circle (\arito);
\draw[fill=white] (Q) circle (\arito);
\draw[fill=white] (-\diametroEntrante,0) circle (\arito);
\draw[fill=white] (0,0) circle (\arito);
\draw[fill=white] (2.0*\kmedios,0) circle (\arito);
\draw[fill=white] (2.0*\kmedios+\diametroEntrante,0) circle (\arito);
\draw (1.4,0.3) node{$\varphi_k$};
\end{scope}
\end{tikzpicture}%
\smallskip
\end{center}

Comprobemos que $\operatorname{long}(PQ) = \varphi_k$ fijándonos en el salino de la derecha.
Por el teorema de Pitágoras, $\operatorname{long}(PO) = \sqrt{1+k^2/4}$; además, $\operatorname{long}(OQ) = k/2$. 
Entonces,
\[
\operatorname{long}(PQ) = \operatorname{long}(PO) + \operatorname{long}(OQ)
= \bigl({k + \sqrt{k^2+4}}\bigr)/{2} = \varphi_k.
\]

%% UNA ALTERNATIVA:
%El diámetro $a$ de los lóbulos laterales también está relacionado con las medias metálicas.
%El ángulo $APD$ es recto, pues abarca el diámetro de la circunferencia,
%con lo cual los triángulos $APB$ y $PDB$ son semejantes.
%% Con $a = \operatorname{long}(AB) = \operatorname{long}(CD)$,
%Ahora, el teorema de Tales nos dice que $a/1 = 1/(k+a)$, luego $a^2+ka-1 = 0$ y, en consecuencia,
%\[
%a = \bigl({-k + \sqrt{k^2+4}}\bigr)/{2} = \varphi_k - k = {1}/{\varphi_k}. 
%% También = -\phi_k 
%\]

% OTRA ALTERNATIVA
El diámetro $a$ de los lóbulos laterales también está relacionado con las medias metálicas.
El ángulo $APD$ es recto, pues abarca el diámetro de la circunferencia,
con lo cual, por el teorema de la altura, $\operatorname{long}(AB) \cdot \operatorname{long}(BD) = 1^2$,
es decir, $a(k+a)= 1$. 
Pero $k+a = \operatorname{long}(BO) + \operatorname{\operatorname{long}}(OD) = \operatorname{long}(OQ) + \operatorname{long}(OP) = \varphi_k$, luego $a = {1}/{\varphi_k}$.

%% Truco para que el t\'{\i}tulo ``Referencias'' no quede tan separado
%% ni destaque tanto el tama\~no de letra:
%\let\section\subsection
%% (se puede hacer pues no hay secciones ni subsecciones;
%% tambi\'en lo suelo usar en la secci\'on de "Problemas y Soluciones")

%%%%%%%%%%%%%%%%%%%%%%%%%%% BIBLIOGRAPHY %%%%%%%%%%%%%%%%%%%%%%%%%%%

%%% No me caben las referencias
%
%\renewcommand{\refname}{Referencia}
%\begin{thebibliography}{1}
%
%\bibitem{C} 
%\textsc{\'O. Ciaurri Ramírez},  El «Libro de los lemas»: un ejercicio de visualización,
%\emph{La Gaceta de la RSME} \textbf{17} (2014), no.~2, 221--245.
%
%\bibitem{P} 
%\textsc{\'A. Plaza}, Relación entre la $k$-ésima media metálica y $\displaystyle \arctan (2/k)$,
%\emph{La Gaceta de la RSME} \textbf{20} (2017), no.~2, 268.
%
%\end{thebibliography}

\medskip
%\vfill

%%%%%%%%%%%%%%%%%%%%%%%%%%%%%%%%%%%%%%%%%%%%%%%%%%%%%%%%%%%%%%%%%%%%%%%
\end{document}
%%%%%%%%%%%%%%%%%%%%%%%%%%%%%%%%%%%%%%%%%%%%%%%%%%%%%%%%%%%%%%%%%%%%%%%

